\documentclass[11pt, a4paper]{article}

% Packages
\usepackage[utf8]{inputenc}
\usepackage[english]{babel}
\usepackage{amsmath}
\usepackage{amssymb}
\usepackage{graphicx}
\usepackage[margin=1in]{geometry}
\usepackage{fancyhdr}
\usepackage{hyperref}
\usepackage{float}
\usepackage{booktabs}
\usepackage{siunitx}

% Header and Footer
\pagestyle{fancy}
\fancyhf{}
\rhead{Measurement Technics - Group 43}
\lhead{KuDAQ Project Report}
\cfoot{\thepage}

% Document Metadata
\title{\textbf{Measurement Technics\\Project Report}\\Group 43 - KuDAQ}
\author{Authors Names}
\date{\today}

\begin{document}

% Title Page
\maketitle

% Table of Contents
\newpage
\tableofcontents
\newpage

% Introduction
\section{Introduction}
This report documents the experimental measurements and analysis conducted in the Measurement Technics course. The project focuses on the development and testing of the KuDAQ measurement system, which integrates various sensors for comprehensive data acquisition and analysis.

\section{Measurements}

\subsection{Measurement 1: Orientation}
\label{sec:measurement1}

\subsubsection{Objective}

The objective of the first measurement is to evaluate the impact of some environmental parameters on the movement of a standardized object caught in a jet flow. In order to do this, we will use the orientation of the object v.s. time, for a given set of parameters, during a fixed time interval. The jet flow is to be generated by a water pump in a water tank of approximately 5L capacity. The standardized object must be small enough to be free to move around the water tank, must not be physically linked to the acquisition system to remain fully independent to avoid potential biais, and must be heavy enough to be able to be caught in the jet flow, but not too heavy to be able to move around the water tank. The parameters that we were able to identify as potentially impacting the movement of the object are the following:

\begin{itemize}
    \item The nozzle shape and size of the water pump
    \item The shape of the object
    \item Some other objects placed in the water tank
    \item The position of the water pump in the water tank
    \item The water level in the water tank
\end{itemize}

We eventually decided for this experiment to focus on the position of the water pump in the water tank, as it is a simple parameter to modify between tests, and that we expect to have a significant impact on the movement of the object. We planned to focus the measurement on the height of the water pump in the water tank, keeping the exact same horizontal position. Because the object will be floating on the water surface, we expect to find larger movement amplitude when the water pump is closer to the water surface, and smaller movement amplitude when the water pump is further from the water surface. 

\subsubsection{Equipment and Setup}

\paragraph{Electrical Hardware}


\paragraph{Mechanical Hardware}
Describe the mechanical hardware and physical setup:
\begin{itemize}
    \item Mechanical mounting and fixtures
    \item Calibration reference devices
    \item Positioning and alignment tools
    \item Environmental control equipment
\end{itemize}

\paragraph{Software Architecture}
Describe the software architecture and computational framework:
\begin{itemize}
    \item Data acquisition software and drivers
    \item Real-time processing framework
    \item Analysis tools (MATLAB/Python)
    \item Communication protocols
\end{itemize}

\paragraph{Acquisition Method}
Describe the data acquisition methodology:
\begin{itemize}
    \item Sampling frequency and resolution
    \item Data storage format
    \item Synchronization method
    \item Trigger and timing specifications
\end{itemize}

\subsubsection{Procedure}
Describe the experimental procedure here. Include:
\begin{enumerate}
    \item Sensor initialization and calibration
    \item Test conditions and environmental parameters
    \item Data collection methodology
    \item Time duration and sampling frequency
\end{enumerate}

\subsubsection{Results}
Present the measurement results and observations:
\begin{itemize}
    \item Raw data from sensors
    \item Processed and filtered data
    \item Graphical representations and plots
    \item Statistical analysis
\end{itemize}

\paragraph{Data Analysis:} Include relevant figures and tables here.

\subsubsection{Discussion}
Discuss the results in the context of:
\begin{itemize}
    \item Expected behavior
    \item Observed deviations and anomalies
    \item Potential sources of error
    \item Accuracy and limitations
\end{itemize}

\subsubsection{Conclusions}
Summarize the key findings from this measurement and their implications.

\newpage

% Additional Measurements
\section{Additional Measurements}

\subsection{Measurement 2: [Title to be determined]}
\label{sec:measurement2}

Add content for additional measurements as needed, following the structure of Measurement 1.

\newpage

% General Conclusion
\section{Conclusion}
Summarize the overall findings of the project, the effectiveness of the KuDAQ system, and recommendations for future work.

\newpage

% References
\begin{thebibliography}{99}
    \bibitem{bmi088} Bosch Sensortec. BMI088 Datasheet.
    \bibitem{madgwick} Madgwick, S. O. H., et al. An efficient orientation filter for inertial and inertial/magnetic sensor arrays.
    % Add more references as needed
\end{thebibliography}

\newpage

% Appendix
\appendix
\section{Raw Data}
Include raw measurement data, extended tables, or additional plots here.

\section{Software and Configuration}
Include code listings, configuration files, or technical documentation here.

\end{document}
